\lecture{4}{Thu 10 Oct 2024 09:32}{Lagrange's Theorem}
\begin{definition}[Index]\label{dfn:aksj}
	Let $G$ be a group with $H\leq G$. Then the index of $H$ in $G$, denoted $[G:H]$, is denoted
	\[
		[G:H]=\frac{\left| G \right| }{\left| H \right| }
	.\]
\end{definition}

Lagrange's theorem is an extremely important result in group theory, and justifies our study of these seemingly arbitrary ``cosets''. One of the immediate corrollaries of Lagrange's theorem generalizes one extremely useful propery of cyclic groups, being that the order of their subgroups divides the order of the group. Many other important theorems also follow from the result, and Lagrange's theorem will remain important throughout our study of algebra.

\begin{theorem}[Lagrange]\label{thm:lagrange}
	Let $G$ be a finite group with subgroup $H\leq G$. Then $[G:H]$ is the number of distinct left cosets of $H$ in $G$. In particular, the order of $H$ divides the order of $G$.
\end{theorem}
\begin{proof}
	Let $\left\{ g_1, \ldots,g_k \right\} \subseteq G$ such that $g_1H, \ldots,g_kH$ are \emph{all of the distinct} cosets of $H$ in $G$. We will show $\left\{ g_1H, \ldots, g_kH \right\} $ is a partition on $G$. By Proposition idk, for any $g_iH,g_jH$ for $i\neq j$, $g_iH\cap g_jH=\varnothing$. Now let $a\in G$, and suppose there does not exist some $g_iH$ such that $a\in g_iH$. Since $H$ is a subgroup, $e\in H$, and thus $ae=a\in aH$. By proposition idk, since  $a\notin g_iH$ for all $i$, we know $aH\neq g_iH$ for all $i$, and thus the set $\left\{ g_1H, \ldots, g_kH \right\} $ is not the set of \emph{all} distinct cosets of $H$ in $G$. By contradiction, and by the methods by which we reached a contradiction, we have $a\in g_iH$ for some $i$, and $\left\{ g_1H, \ldots, g_kH \right\} $ partitions $G$.

	By Proposition idk, note $\left| g_iH \right| =\left| g_jH \right| $ for all $i,j$. Since $g_1H, \ldots, g_kH$ forms a parition of $G$, we have
	\[
		\left| G \right| =\sum_{i=1}^{k} \left| g_iH \right| =k\left| g_iH \right| 
	.\]
	Therefore,
	\[
		\frac{\left| G \right| }{\left| g_iH \right| }=k
	,\]
	and $k$ is the number of cosets. However, there exists some $g_i=e$, and $eH=H$, so we have
	\[
		\frac{\left| G \right| }{\left| H \right| }=k \implies [G:H]=k
	.\]
\end{proof}
\begin{remark}
	didnt judson define $[G:H]$ the number of cosets and show
	\[
		[G:H]=\frac{\left| G \right| }{\left| H \right| }
	?\]
\end{remark}
Now we will explore the many corrolaries of \ref{thm:lagrange}.
\begin{corollary}\label{cor:as}
	Let $G$ be a finite group. If $g\in G$, then $\left| g \right| $ divides $\left| G \right| $.
\end{corollary}
\begin{proof}
	Let $H=\left<g \right>$. Thus, $\left| H \right| =\left| g \right| $. By \ref{thm:lagrange}, $\left| H \right| $ divides $\left| G \right| $, so $\left| g \right| $ divides $\left| G \right| $.
\end{proof}
\begin{corollary}
	Let $G$ be a finite group. For all $a\in G$, $a^{\left| G \right| }=e$.
\end{corollary}
\begin{proof}
	By \ref{cor:as}, for any $a\in G$, $\left| a \right| |\left| G \right| $. Thus,
	\[
		\left| G \right| =k\left| a \right| 
	\]
	for some $k\in\mathbb{Z}$, and thus
	\[
		\left| a \right| =\frac{\left| G \right| }{k}\in\mathbb{N}
	.\]
	We then have
	\[
		a^{\left| G \right| /k}=e
	.\]
	Thus,
	\[
		a^{\left| G \right| }=a^{k (\left| G \right| /k)}=e^k=e	
	.\]
\end{proof}
\begin{definition}[Prime Number]\label{dfn:bwa}
	A positive number $p\geq 2$ is prime if and ony if the positive divisors of $p$ are precisely $1$ and $p$.
\end{definition}
\begin{proposition}
	Let $p\in\mathbb{N}$. Then $p$ is prime if and only if when $p|ab$, then $p|a$ or $p|b$.
\end{proposition}
\begin{corollary}\label{coror}
	Let $\left| G \right| =p$ for a prime number $p$ prime. Then $G$ is cyclic, generated by any $a\in G\setminus\left\{ e \right\} $.
\end{corollary}
\begin{proof}
	Let $a\in G\setminus\left\{ e \right\} $. Let $H=\left<a \right>$. It follows that $\left| H \right| \geq 2$, since $a\neq e$, so $a,e\in H$. By \ref{thm:lagrange}, the order of $H$ divides $\left| G \right| $. Since $\left| G \right| $ is prime, by \ref{dfn:bwa}, the only divisors of $\left| G \right| $ are $1$ and $\left| G \right|$. Since $\left| H \right| >1$, we have $\left| H \right| =\left| G \right| $, and thus $H=G$.
\end{proof}
\begin{corollary}
	Let $G$ be a finite group with $\left| G \right| =p$ prime. Then $G\cong \mathbb{Z}_p$.
\end{corollary}
\begin{proof}
	By \ref{coror}, $G$ is cyclic with order $p$. By an earlier theorem I still have to prove, $G\cong \mathbb{Z}_{\left| G \right| }=\mathbb{Z}_p$.
\end{proof}
\begin{corollary}
	If $H$ and $K$ are subgroups of a finite group $G$ with $K\subset H\subset G$, then
	\[
		[G:K]=[G:H][H:K]
	.\]
\end{corollary}
\begin{proof}
	\[
		[G:K]=\frac{\left| G \right| }{\left| K \right| }=\frac{\left| G \right| }{\left| H \right| }\cdot \frac{\left| H \right| }{\left| K \right| }
	.\]
\end{proof}
\begin{remark}
	The converse of Lagrange's Theorem is \emph{not true}. That is, if $b|\left| G \right| $, is possible for there to be no subgroups of order $b$.
\end{remark}
\begin{theorem}[Fermat's Little Theorem]\label{thm:fermat}
	Let $p$ be a prime number and suppose $p$ does not divide $a$. Then
	\[
		a^{p-1}\equiv 1\mod{p}
	.\]
	Equivalently,
	\[
		a^p=a\mod{p}
	.\]
\end{theorem}
\begin{proof}
	Since $a$ does not divide $p$, we know $a=pm+r$ for $0<r<p$. We then know
	\[
	a^{p-1} \equiv \left( pm+r \right) ^p\equiv  r^{p-1}\mod{p}
	\]
	by like 10 different theorems, such as the binomial theorem. We can then assume $r\in U(p)$, since $r\neq p$. Since $\left| U(p) \right| =p-1$ then $r^{p-1}\equiv 1\mod{p}$.
\end{proof}
\begin{theorem}[Generalized Fermat]\label{thm:lask}
	Let $a$ be relatively prime to $n$. Then $a^{\phi (n)}\equiv 1\mod{n}$.
\end{theorem}
\begin{proof}
	Since $a$ is relatively prime to $ n$, then $a\in U(n)$. Note that $\phi (n)=\left| U(n) \right| $. Thus, $a^{\phi (n)}=a^{\left| U(n) \right| }$. By \ref{cor:as}, $\left| a \right| $ divides $\left| U(n) \right| $, and thus $a^{\left| U(n) \right| }\equiv 1\mod{n}$.
\end{proof}
\section{Application to Permutation Groups}
This is a window into ``group actions''. The book does not cover group actions, so I'll ask the prof for resources about it.
\begin{definition}[Stabilizer and Orbit]\label{dfn:lk}
	Let $G$ be a group of permutations of some set $S$. Then for each $i\in S$, define the \emph{stabilizer} of $i$ in $G$,
	\[
		\operatorname{stab}_G(i) =\left\{ \phi \in G \mid \phi (i)=i \right\} 
	,\]
	and the \emph{orbit} of $i$ in $G$,
	\[
		\operatorname{orb}_G(i) =\left\{ \phi (i) \mid \phi \in G \right\} 
	.\]
\end{definition}
\begin{theorem}\label{thm:ksnd}
	$\operatorname{stab}_G(i) \leq G$. $\operatorname{orb}_G(i) \subseteq S$.
\end{theorem}
\begin{proof}
	Let $G$ be a group of permutations of $S$, and let $i\in S$. First, we show $\operatorname{stab}_G(i) $ is closed under function composition. Let $\phi ,\psi \in \operatorname{stab}_G(i) $. Since all $\phi \in \operatorname{stab}_G(i) $, have $\phi (i)=i$, it follows that
	\[
		\left( \phi \psi  \right) (i)=\phi (\psi (i))=\phi (i)=i
	.\]
	Therefore, $\phi \psi \in \operatorname{stab}_G(i) $.

	Associativity follows from the fact that $\operatorname{stab}_G(i) \subseteq G$, which is a group. Notice that the identity map $\mathbbm{1}_S$ maps $i\mapsto i$, so $\mathbbm{1}_S\in \operatorname{stab}_G(i) $. Finally, we show inverses exist. Let $\phi \in \operatorname{stab}_G(i) $, and let $\phi ^{-1}$ be its inverse in $G$. It follows that
	\[
		\left( \phi ^{-1}\phi  \right) (i)=i
	.\]
	However, $\phi (i)=i$, so it follows that
	\[
		\phi ^{-1}(\phi (i))=\phi ^{-1}(i)=i
	.\]
	Therefore, $\phi ^{-1}$ fixes $i$, and thus $\phi ^{-1}\in \operatorname{stab}_G(i) $. Thus, $\operatorname{stab}_G(i) $ is a group. Since $\operatorname{stab}_G(i) \subseteq G$, we have $\operatorname{stab}_G(i) \leq G$.

	The second fact is trivial. For any $\phi \in G$, we know $\phi $ is by definition a bijection $S\hookrightarrow S$, so $\phi (i)\in S$ for all $i\in S$. Therefore, any $i'\in \operatorname{orb}_G(i) $ has $i'=\phi (i)$ for some $\phi \in G$, and thus $i'\in S$. Therefore, $\operatorname{orb}_G(i) \subseteq S$.
\end{proof}
\begin{theorem}[Orbit-Stabilizer]\label{thm:a}
	Let $G$ be a finite group of permutations of a set $S$. Then for any  $i\in S$,
	\[
		\left| G \right| =\left| \operatorname{orb}_G(i)  \right| \left| \operatorname{stab}_G(i)  \right| 
	.\]
\end{theorem}
\begin{proof}
	this one is harder
\end{proof}
\chapter{External Direct Products}
